\documentclass[aspectratio=169]{beamer}

\usetheme{colorful-dream}

\usepackage{minted}
\setminted{fontsize=\footnotesize, breaklines=true}

\title{Podstawy Pythona dla Data Science}
\subtitle{Wykład 2: NumPy, Matplotlib i sygnały inżynierskie}
\author{Lu Pa}
\date{\today}

% --- Title Slide Logos ---
\titlegraphic{
  \vspace{0.5cm}
  \begin{center}
    \includegraphics[height=1.8cm]{logo_university_placeholder.png}
    \hspace{1cm}
    \includegraphics[height=1.8cm]{logo_faculty_placeholder.png}
  \end{center}
}

\begin{document}

% ---------------------------------------------------------
\begin{frame}
  \titlepage
\end{frame}

% ---------------------------------------------------------
\section{PART I}
\begin{frame}[plain]
  \begin{center}
    \vspace{2cm}
    {\usebeamerfont{title}\color{white}\Huge PART I}
  \end{center}
\end{frame}

% ---------------------------------------------------------
\begin{frame}{Od Pythona do obliczeń naukowych}
  \begin{itemize}
    \item Python staje się potężny dzięki bibliotekom naukowym.
    \item Dziś skupiamy się na:
      \begin{itemize}
        \item \textbf{NumPy} — tablice numeryczne i wektoryzacja
        \item \textbf{Matplotlib} — wykresy inżynierskie
      \end{itemize}
    \item Te dwie biblioteki są fundamentem ekosystemu naukowego.
  \end{itemize}
\end{frame}

% ---------------------------------------------------------
\begin{frame}{Czym jest NumPy?}
  \begin{itemize}
    \item Podstawowa biblioteka do obliczeń numerycznych.
    \item Zapewnia:
      \begin{itemize}
        \item szybkie tablice wielowymiarowe
        \item operacje wektorowe
        \item funkcje matematyczne
        \item generatory liczb losowych
      \end{itemize}
    \item Napisana w C dla wysokiej wydajności.
    \item Podstawa Pandas, SciPy, scikit-learn i wielu innych.
  \end{itemize}
\end{frame}

% ---------------------------------------------------------
\begin{frame}[fragile]{Tworzenie tablic NumPy}
\begin{minted}{python}
import numpy as np

a = np.array([1, 2, 3])
b = np.linspace(0, 10, 100)
c = np.zeros(50)
\end{minted}
\end{frame}

% ---------------------------------------------------------
\begin{frame}[fragile]{Operacje wektorowe}
\begin{minted}{python}
v = np.array([10, 20, 30])
v_ms = v / 3.6
\end{minted}

\begin{itemize}
  \item Operacje działają na całych tablicach.
  \item Dużo szybsze i czytelniejsze niż pętle.
\end{itemize}
\end{frame}

% ---------------------------------------------------------
\begin{frame}{Przykład inżynierski: sygnał prędkości}
  \begin{itemize}
    \item Wyciągamy tablice:
      \begin{itemize}
        \item czas
        \item prędkość
      \end{itemize}
    \item Obliczamy:
      \begin{itemize}
        \item przyspieszenie
        \item jerk
        \item wygładzoną prędkość
      \end{itemize}
  \end{itemize}
\end{frame}

% ---------------------------------------------------------
\begin{frame}[fragile]{Wyciąganie tablic z DataFrame}
\begin{minted}{python}
t = df["time_s"].values
v = df["speed_kmh"].values
\end{minted}
\end{frame}

% ---------------------------------------------------------
\begin{frame}[fragile]{Przyspieszenie — różniczkowanie numeryczne}

\begin{equation}
a(t) \approx \frac{\Delta v}{\Delta t}
\end{equation}

\begin{minted}{python}
dt = np.diff(t)
dv = np.diff(v)
accel = dv / dt
\end{minted}
\end{frame}

% ---------------------------------------------------------
\begin{frame}[fragile]{Jerk — pochodna przyspieszenia}

\begin{equation}
j(t) \approx \frac{\Delta a}{\Delta t}
\end{equation}

\begin{minted}{python}
jerk = np.diff(accel) / dt[:-1]
\end{minted}
\end{frame}

% ---------------------------------------------------------
\begin{frame}[fragile]{Wygładzanie sygnału — średnia krocząca}
\begin{minted}{python}
window = 15
v_smooth = np.convolve(v, np.ones(window)/window, mode="same")
\end{minted}
\end{frame}

% ---------------------------------------------------------
\begin{frame}{Matplotlib — wykresy inżynierskie}
  \begin{itemize}
    \item Elastyczna i potężna biblioteka do wizualizacji.
    \item Idealna do:
      \begin{itemize}
        \item przebiegów czasowych
        \item dashboardów wielopanelowych
        \item histogramów
        \item wykresów punktowych
      \end{itemize}
    \item Bardzo konfigurowalna.
  \end{itemize}
\end{frame}

% ---------------------------------------------------------
\section{PART II}
\begin{frame}[plain]
  \begin{center}
    \vspace{2cm}
    {\usebeamerfont{title}\color{white}\Huge PART II}
  \end{center}
\end{frame}

% ---------------------------------------------------------
\begin{frame}[fragile]{Wielopanelowy wykres inżynierski}
\begin{minted}{python}
fig, ax = plt.subplots(2, 2, figsize=(12, 8))

ax[0,0].plot(t, v)
ax[0,0].set_title("Prędkość [km/h]")

ax[0,1].plot(t[:-1], accel)
ax[0,1].set_title("Przyspieszenie [m/s²]")

ax[1,0].plot(t, v_smooth)
ax[1,0].set_title("Wygładzona prędkość")

ax[1,1].plot(t, df["engine_rpm"])
ax[1,1].set_title("Obroty silnika")

for a in ax.ravel():
    a.grid(True)

plt.tight_layout()
plt.show()
\end{minted}
\end{frame}

% ---------------------------------------------------------
\begin{frame}{Zadania praktyczne}
  \begin{itemize}
    \item Oblicz jerk i go zwizualizuj.
    \item Porównaj różniczkowanie numeryczne z \texttt{np.gradient}.
    \item Przetestuj różne okna wygładzania.
    \item Wykonaj histogram przyspieszenia.
    \item Zbuduj dashboard 2×2 z kluczowymi sygnałami.
  \end{itemize}
\end{frame}

% ---------------------------------------------------------
\end{document}
